% Ссылка в тексте: (приложение~И)
\chapter{ОЦЕНОЧНЫЕ АРТЕФАКТЫ И ТЕКУЩИЕ РЕЗУЛЬТАТЫ}

В~приложении собраны те артефакты валидации, которые уже подготовлены
или~измерены в~текущей редакции диплома.

\section*{И.1 --- Агрегированные результаты конфигурации B2}

\begin{table}[h!]
  \caption{Агрегированные результаты пилотного прогона одношаговой конфигурации B2}
  \label{tab:appendix-b2}
  \centering
  \begin{tabular}{|p{4.5cm}|p{2.0cm}|p{2.2cm}|p{2.5cm}|}
  \hline
  \textbf{Срез корпуса} & \textbf{$n$} & \textbf{Покрытие, \%} & \textbf{Доля $\geq 70\%$} \\
  \hline
  Все документы & 19 & 61.4 & 21.1\% \\
  \hline
  Продуктово-редакционные документы & 7 & 73.8 & 28.6\% \\
  \hline
  Корпоративные регламенты & 6 & 38.9 & 0\% \\
  \hline
  Учебные программы MIT~OCW & 6 & 69.5 & 33.3\% \\
  \hline
  \end{tabular}
\end{table}

Эти значения относятся к~baseline B2 из артефакта
\texttt{baseline\_comparison.json} и используются как технический
ориентир MVP. Они не подменяют paired B2/B4-сравнение, которое должно
быть выполнено на том же корпусе после стабильного B4-прогона.

\section*{И.2 --- Подготовленные оценочные артефакты}

\begin{table}[h!]
  \caption{Оценочные наборы текущего MVP}
  \label{tab:appendix-assets}
  \centering
  \begin{tabular}{|p{5.0cm}|p{2.0cm}|p{5.5cm}|}
  \hline
  \textbf{Артефакт} & \textbf{Объём} & \textbf{Назначение} \\
  \hline
  Измеренный корпус baseline B2 & 19 & подготовка курсов и~сравнение базовых подходов \\
  \hline
  Вопросно-ответные пары (\texttt{qa\_pairs.json}) & 30 & оценка faithfulness и answer relevancy \\
  \hline
  Фактуальные утверждения (\texttt{statements.json}) & 100 & расчёт precision / recall / F1 для~факт-чекинга \\
  \hline
  \end{tabular}
\end{table}

\section*{И.3 --- Экспертная и пользовательская валидация: статус}

На~момент подготовки текущего черновика:
\begin{itemize}
  \item шаблон слепой экспертной оценки подготовлен, но таблицы B1--B4
    ещё не~заполнены измеренными значениями;
  \item пилотный протокол с~метриками SUS / NPS сформирован, но не~представляется
    как завершённый численный результат;
  \item приложение фиксирует именно \textit{готовность контуров оценки},
    а~не выдаёт незавершённые измерения за~готовый оценочный результат.
\end{itemize}

\section*{И.4 --- Машинно читаемые evidence-артефакты платформы}

В текущей кодовой базе ИИ-методиста уже предусмотрены поля, необходимые
для дипломной валидации:
\begin{itemize}
  \item \texttt{citations} и \texttt{block\_provenance} --- ссылочная
    привязка уроков к источникам;
  \item \texttt{quality\_seal} --- статус \texttt{PASS},
    \texttt{PASS\_WITH\_NOTES} или \texttt{FAIL};
  \item \texttt{cost\_quality} --- стоимость, модель, токены и
    объяснение quality trade-offs;
  \item \texttt{thesis\_validation} --- snapshot метрик H1--H3,
    пригодный для сохранения в экспериментальный корпус.
\end{itemize}

Фронтовый E2E-тест \texttt{ai-methodist-thesis-validation.spec.ts}
проверяет отображение этих полей в пользовательском контуре, а скрипт
\texttt{run\_ai\_methodist\_product\_validation.py} превращает JSON
артефакты в gate-report по гипотезам диплома.
